%---------------------------------------------------------
\subsection{Vista A -- Responsabile delle operazioni di magazzino}
\label{sec:req-A}
Il responsabile delle operazioni di magazzino gestisce i flussi fisici della merce, ed è un dipendente del committente che lavora con il Warehouse Management System del software SYSTORE. Registra l’arrivo delle unità di carico dalla produzione, consulta le giacenze di ogni articolo e ubicazione, blocca e sblocca i lotti (individuando quelli con la scadenza più vicina), rilascia gli ordini di spedizione e controlla le missioni in corso su ogni macchina.
Per compiere ciò, per ogni articolo si registrano codice, descrizione, formato di vendita, elementi per pallet e durata di conservazione. La merce arriva in lotti di produzione: ogni lotto è la controparte fisica di un articolo ed è identificato dal codice dell’articolo stesso e da un secondo identificatore numerico, univoco all’interno dello stesso prodotto. Presenta inoltre la data di produzione e di scadenza. 
La merce è spostata e sistemata su unità di carico, le quali si dividono in pallet e contenitori, questi ultimi presenti solo nelle aree di picking. Ogni unità ha un codice univoco, un tipo, un’altezza e un peso misurati al ricevimento. Solitamente le unità contengono un solo tipo di articolo, ma in fase di picking è possibile comporre un pallet inserendone più di uno; per ogni articolo selezionato si registrano lotto e numero di elementi.
Un’ubicazione è una posizione che può ospitare una unità di carico, con un codice univoco e uno stato (libera, occupata o bloccata). Le ubicazioni si dividono in celle e baie. Il magazzino è diviso in corridoi, ognuno servito da un traslo-elevatore, e ogni corridoio è composto da celle: una cella è individuata nel corridoio da lato, colonna, livello e profondità, ha altezza e portata massime e ospita al più una unità di carico. Le baie sono i punti di scambio con l’esterno del magazzino e hanno un tipo, che ne specifica il compito rispetto al magazzino (come ingresso, uscita, picking o carico). 
Ogni spostamento di una unità da un’ubicazione di origine a una di destinazione è una missione, creata dal WMS con codice, priorità, stato (in coda, in esecuzione, completata o annullata) e istanti di creazione, inizio e fine. Le missioni si distinguono in missioni di ingresso, di uscita e di trasferimento: le prime portano una UDC da una baia a una cella e registrano l’esito del controllo di sagoma e peso, per assicurarsi la correttezza del prodotto; le seconde la portano da una cella a una baia per evadere un ordine; le terze la spostano tra due celle o tra due baie e ne registrano il motivo. Una missione impiega di solito più macchine in sequenza e per ognuna si registrano la posizione nella sequenza e gli istanti di inizio e fine del tratto.
Le missioni sono eseguite dalle macchine dell’impianto: ognuna di esse può essere la rappresentazione di un modello scelto tra traslo-elevatori, trasportatori e veicoli laser comandati. Di ogni macchina si consultano codice, descrizione e stato (in lavoro, in attesa, in errore, in manutenzione). 
Gli ordini di spedizione arrivano dal gestionale del committente verso SYSTORE, ciascuno con numero, cliente destinatario, data di spedizione prevista, baia di carico e stato, e con il numero di elementi richiesto per ogni articolo. 
